% !TeX root = main.tex
% !TeX program = xelatex
% Introduction follows the analysis scope and evidence in Chapters 2--5.

\section{引言}
\label{sec:introduction}

代理间安全通信需要同时验证参与方身份、判定联系资格并约束后续请求次数。
SAGA 通过接触策略、一次性公钥（One-Time Key，OTK）和访问控制令牌（Access Control Token，ACT）实现分阶段通信授权~\cite{syros2025sagaV2}。
Provider 根据接触策略与领取额度发放 OTK，接收方执行本地策略检查并消费 OTK，为实际连接者签发 ACT，后续请求由 ACT 的剩余额度控制。
规则判定、凭证使用与额度消耗共同决定这一机制所能提供的授权保证。

分阶段授权将策略决策、凭证消费与请求计数分布在不同主体和状态中。
已有授权策略分析能够检查规则组合是否产生非预期许可，以及策略变更如何影响访问结果~\cite{fisler2005margrave}。
在 SAGA 中，Provider 按领取身份维护 OTK 发放额度，接收方依据实际连接身份、本地策略及 OTK 未消费状态签发 ACT，并在多个连接之间共享同一令牌的剩余额度~\cite{syros2025sagaV2,sagaReferenceImpl2026}。
对这一机制的分析需要确定实际规则选择是否符合预期优先关系、按身份维护的领取额度是否足以限制同一身份取得新 ACT 的次数，以及并发检查与扣减能否共同维持累计准入上限。

围绕这些执行约束，采用参考语义对照、符号协议分析与共享状态模型检查相结合的方法。
策略分析独立计算预期预算，将其与原生匹配函数对照，并通过受控消费者执行追踪决策对 ACT 建立的影响。
OTK 分析采用 Tamarin~\cite{meier2013tamarin} 追踪发放、主动转交与 ACT 建立中的身份关系，检查成功使用的 OTK 是否曾发放给同一身份，同时保留发放上限和单次消费性质。
ACT 分析以 TLA$^{+}$~\cite{lamport2002specifying} 表示多个连接访问共享额度的执行交错，通过 TLC 检查累计准入不变量，并以受控原生方法实验观察后端调用次数。

分析揭示了三个阶段的约束执行问题。
策略匹配器未正确维护规则优先关系，受控执行中预期拒绝的发送方取得并解密了 ACT。
OTK 符号反例中，两端策略均允许的发送方可在耗尽自身领取额度后，利用另一身份主动转交的未消费 OTK 建立并解密新 ACT。
该配置中按身份发放上限与 OTK 单次消费均成立，表明发放计数与按实际使用身份限制新 ACT 次数之间存在差异。
ACT 的检查与扣减分离使 TLC 检查出现超额准入反例，原生方法级受控实验进一步观察到额外后端调用。
这些结果将 SAGA 的授权执行缺口定位于规则选择、领取与使用身份的关联以及共享额度的并发更新。
